\begin{table}[t]
  \centering
  \caption{\textbf{合成 test を撮影範囲の被覆率で層別した内訳（ours）。}
  層は GT のコート線サンプルのうち画像内に落ちる割合で分けた。
  \texttt{full} はコート全体、\texttt{near\_full} はほぼ全体、
  \texttt{partial} は一部だけが画面内に入る。
  coverage の低い層ほど誤差の裾が伸びるが、
  どの層でも PCK@0.01 は 0.92 を下回らない。
  }
  \label{tab:comparison_strata}
  \small
  \begin{tabular}{@{}lrrrr@{}}
    \toprule
    層 & \(n\) & 画面内線 median (px) & PCK@0.01 & GT-visible 点 \\
    \midrule
    \texttt{full} & \resSynthOursFullN &
      \resSynthOursFullLineMedian & \resSynthOursFullPckOne & \resSynthOursFullVisible \\
    \texttt{near\_full} & \resSynthOursNearFullN &
      \resSynthOursNearFullLineMedian & \resSynthOursNearFullPckOne & \resSynthOursNearFullVisible \\
    \texttt{partial} & \resSynthOursPartialN &
      \resSynthOursPartialLineMedian & \resSynthOursPartialPckOne & \resSynthOursPartialVisible \\
    \bottomrule
  \end{tabular}
\end{table}

\begin{table}[t]
  \centering
  \caption{\textbf{合成 test の会場別 PCK@0.01（ours）。}
  test は trajectory-group-disjoint であり、
  会場自体は学習にも現れる。したがってこの表は
  未知会場への汎化ではなく、
  同じ会場の未学習軌道でのばらつきを示す。
  }
  \label{tab:comparison_scene}
  \small
  \begin{tabular}{@{}lrr@{}}
    \toprule
    会場 & \(n\) & PCK@0.01 \\
    \midrule
    \texttt{B00} & \resSynthOursSceneBzeroN & \resSynthOursSceneBzeroPckOne \\
    \texttt{B01} & \resSynthOursSceneBoneN & \resSynthOursSceneBonePckOne \\
    \texttt{B02} & \resSynthOursSceneBtwoN & \resSynthOursSceneBtwoPckOne \\
    \texttt{B03} & \resSynthOursSceneBthreeN & \resSynthOursSceneBthreePckOne \\
    \bottomrule
  \end{tabular}
\end{table}
